
\slajdid{05-s029}
\begin{frame}[label=05-s029]
\gusttekst
\begin{center}Prekoračenje opsega –zašto nam treba informacija?\end{center}
\begin{columns}[c,onlytextwidth]
\column{.26\textwidth}Zbir dve negativne vrednosti u slučaju OVF daje pozitivnu vrednost
\column{.40\textwidth}\centering\def\DEsaturacija{0}\begin{tikzpicture}[x=.8cm,y=.8cm,font=\fontsize{9}{10}\selectfont]
\draw[DEplava,->] (-2.3,0)--(2.5,0)node[below]{$a+b$};
\draw[DEplava,->] (0,-1.25)--(0,1.45)node[left,text=DEtekst]{$s$};
\ifnum\DEsaturacija=0
 \draw[DEplava] (-2.3,-.8)--(-.8,.8)--(-.8,-.8)--(.8,.8)--(.8,-.8)--(2.3,.8);
 \draw[orange,->] (-1.1,0)--(-1.1,.45);\draw[orange,->] (-1.1,.45)--(0,.45);
 \draw[orange,->] (1.2,0)--(1.2,-.4);\draw[orange,->] (1.2,-.4)--(0,-.4);
 \draw[DEzelena] (-.8,-.8)--(-.8,-1.4) (.8,-.8)--(.8,-1.4);
 \draw[DEzelena,<->] (-.8,-1.3)--(.8,-1.3)node[midway,below,text=DEtekst]{regularno};
\else
 \draw[DEplava] (-2.1,-.8)--(-.8,-.8)--(.8,.8)--(1.9,.8);
 \draw[orange,->] (-1.1,0)--(-1.1,-.8);\draw[orange,->] (-1.1,-.8)--(0,-.8);
 \draw[orange,->] (1.2,0)--(1.2,.8);\draw[orange,->] (1.2,.8)--(0,.8);
\fi
\end{tikzpicture}

\column{.30\textwidth}Zbir dve pozitivne vrednosti u slučaju OVF daje negativnu vrednost
\end{columns}\par\vspace{4mm}
\begin{columns}[c,onlytextwidth]
\column{.26\textwidth}
\column{.40\textwidth}\centering\def\DEsaturacija{1}\begin{tikzpicture}[x=.8cm,y=.8cm,font=\fontsize{9}{10}\selectfont]
\draw[DEplava,->] (-2.3,0)--(2.5,0)node[below]{$a+b$};
\draw[DEplava,->] (0,-1.25)--(0,1.45)node[left,text=DEtekst]{$s$};
\ifnum\DEsaturacija=0
 \draw[DEplava] (-2.3,-.8)--(-.8,.8)--(-.8,-.8)--(.8,.8)--(.8,-.8)--(2.3,.8);
 \draw[orange,->] (-1.1,0)--(-1.1,.45);\draw[orange,->] (-1.1,.45)--(0,.45);
 \draw[orange,->] (1.2,0)--(1.2,-.4);\draw[orange,->] (1.2,-.4)--(0,-.4);
 \draw[DEzelena] (-.8,-.8)--(-.8,-1.4) (.8,-.8)--(.8,-1.4);
 \draw[DEzelena,<->] (-.8,-1.3)--(.8,-1.3)node[midway,below,text=DEtekst]{regularno};
\else
 \draw[DEplava] (-2.1,-.8)--(-.8,-.8)--(.8,.8)--(1.9,.8);
 \draw[orange,->] (-1.1,0)--(-1.1,-.8);\draw[orange,->] (-1.1,-.8)--(0,-.8);
 \draw[orange,->] (1.2,0)--(1.2,.8);\draw[orange,->] (1.2,.8)--(0,.8);
\fi
\end{tikzpicture}

\column{.30\textwidth}\centering Česta realizacija u digitalnoj obradi signala\par\vspace{4mm}U slučaju OVF zasićenje
\end{columns}
\end{frame}
